\setwlchapter{correlated-solutions}

\section*{Correlated electrons}
The checks below form a fresh session. In each case the calculation uses
a different representation or route from the analytical argument.

\begin{solution}{ex:defects}
At fixed $N=L$, a state with no doublons has one electron on every
site. The spins can be chosen independently, giving $2^L$ ground
states. A state at energy $U$ has precisely one doublon and one hole.
Choose the doublon site in $L$ ways and the hole site in $L-1$ ways.
There is no spin choice on either of these two sites, giving
$L(L-1)2^{L-2}$ states. For four sites the two counts are 16 and 48.
The other six half-filled states have two doublons and two holes.
Their inclusion gives $\binom84=70$ as required. With $U$ as the
energy unit the full half-filled matrix is diagonal:
\begin{wl}
Needs["sneg`"];
snegfermionoperators[c, d];
sites4 = c /@ Range[4];
b4 = Select[qbasisvc[sites4], First[#] == {0} &];
atomic4 = First[makematricesbzvc[Total[hubbard[sites4]], b4]][[2]];
TableForm[Sort[Tally[Eigenvalues[atomic4]]],
  TableHeadings -> {None, {"Energy / U", "States"}}]
\end{wl}
\end{solution}

\begin{solution}{ex:detuned}
The singly occupied configurations have zero site energy. The two
intermediate doublons cost $U-\delta$ and $U+\delta$. Each is reached
from the singlet with squared amplitude $2t^2$, while the triplet
cannot hop. Consequently
\[
 J(\delta)=2t^2\left(\frac1{U-\delta}+\frac1{U+\delta}\right)
 =\frac{4t^2 U}{U^2-\delta^2}.
\]
Positive excitation costs require $|\delta|<U$, and we need
$|t|\ll U-|\delta|$. The divergence as $|\delta|\to U$ signals
that a charge state must be retained in the low-energy subspace.
It is not a physical divergence of the exchange energy. Verify at
finite detuning by diagonalizing the full singlet matrix; the two
answers are printed in the same energy units:
\begin{wl}
snegrealconstants[t, U, delta];
sites2 = {c[1], c[2]};
hd = U Total[hubbard[sites2]] - t hop[c[1], c[2]];
biased = hd + delta (number[c[2]] - number[c[1]])/2;
bs = Select[qsbasisvc[sites2], First[#] == {0, 0} &];
hs = First[makematricesbzvc[biased, bs]][[2]];
par = {U -> 10, delta -> 6, t -> 1/10};
j2 = 4 t^2 U/(U^2 - delta^2);
N[{-Min[Eigenvalues[hs /. par]], j2 /. par}, 6]
\end{wl}
The triplet energy is exactly zero even at finite detuning. The
difference between the two printed splittings is the contribution of
fourth and higher even orders.
\end{solution}

\begin{solution}{ex:gaps}
One electron occupies a bonding orbital of energy $-t$. A three-electron
state contains exactly one doublon, so its lowest energy is $U-t$.
The half-filled ground-state energy is $E_-$ of Eq.~\eqref{eq:dimer}.
Substitution in Eq.~\eqref{eq:chargegap} gives
\[
 \Delta_c=\sqrt{U^2+16t^2}-2t,\qquad
 \Delta_s=(\sqrt{U^2+16t^2}-U)/2.
\]
For $U=0$ both are $2t$: the bonding level forms a closed shell, and
the antibonding level lies $2t$ higher. At $U\gg t$,
$\Delta_c=U-2t+8t^2/U+\cdots$, while $\Delta_s=4t^2/U+\cdots$.
As a direct check, set $(t,U)=(1,4)$, use charge \emph{only}, and
take the lowest eigenvalue separately in each particle-number sector:
\begin{wl}
bc = qbasisvc[sites2];
mc = makematricesbzvc[hd /. {t -> 1, U -> 4}, bc];
e[q_] := Min[Eigenvalues[First[Select[mc, First[#] == {q} &]][[2]]]];
Simplify[e[1] + e[-1] - 2 e[0] == Sqrt[32] - 2]
\end{wl}
\end{solution}

\begin{solution}{ex:moments}
Expand $G_\sigma$ in inverse powers of $z$:
\[
 G_\sigma(z)=\sum_{\ell\ge0}M_\ell/z^{\ell+1}.
\]
Equation~\eqref{eq:atomicG} then gives
$M_\ell=(1-w)\epsilon^\ell+w(\epsilon+U)^\ell$. Thus
\[
 M_0=1,\quad M_1=\epsilon+Uw,\quad
 M_2=\epsilon^2+(2\epsilon U+U^2)w.
\]
A single pole at the Hartree energy reproduces $M_0$ and $M_1$,
but has $M_2=M_1^2$. The exact variance $M_2-M_1^2=U^2w(1-w)$
vanishes only when the peaks coincide or one carries zero weight.
For an independent computation \fn{spectralmoments} evaluates
$M_\ell=\langle\{ \mathcal{L}^\ell d_\uparrow,d^\dagger_\uparrow\}\rangle$
with $\mathcal{L}A=[A,H]$:
\begin{wl}
snegrealconstants[\[Epsilon]];
ha = \[Epsilon] number[d[]] + U hubbard[d[]];
mom = spectralmoments[d[AN, UP], d[CR, UP], ha, 2] /.
  avg[number[d[], DO]] -> w;
Expand[mom - {1, \[Epsilon] + U w, \[Epsilon]^2 + (2 \[Epsilon] U + U^2) w}]
\end{wl}
\end{solution}
